\شروع‌مسئله[50] \متن‌سیاه{پیشنهاد الگوریتم یا داده‌ساختار مناسب}. در هر قسمت از این سؤال از شما خواسته می‌شود یک الگوریتم یا داده‌ساختار مناسب پیشنهاد دهید. این داده‌ساختار یا الگوریتم می‌تواند از مطالب گفته شده در درس، ترکیبی از آن‌ها و یا خلاقیتی باشد. علاوه بر \زیرخط{کلیت داده‌ساختار و الگوریتم}، \زیرخط{زمان اجرا} و \زیرخط{دلیل انتخاب خود} را نیز در ۲ الی ۳ خط (قسمت‌های آ تا ج) و ۴ الی ۵ خط (قسمت‌های چ تا ذ) به طور واضح توضیح دهید. در این سؤال نیاز به اثبات درستی الگوریتم و اثبات زمان اجرای الگوریتم نیست. همچنین در صورتی که زمان اجرای الگوریتم شما \زیرخط{زمان سرشکن یا زمان تصادفی} است، آن را بیان کنید.
% صورت اصلی مسئله را در این‌جا بنویسید.
\شروع{زیرمسائل}
\زیرمسئله[4]
یک آرایه در ابتدا مرتب بوده، ولی ۱۰ عنصر آن (که نمی‌دانیم کدام عناصر هستند) توسط یک نفر تغییر داده شده است. چه الگوریتمی برای مرتب‌سازی این آرایه پیشنهاد می‌کنید؟

\پاسخ{
مرتب‌سازی درجی. این مرتب‌سازی برای آرایه‌هایی که تقریباً مرتب هستند مناسب است و این مسئله را در $O(n)$ حل می‌کند.
	
}
\نمره‌دهی{
	کلیت الگوریتم ۲ نمره و زمان اجرای آن نیز ۲ نمره دارد.
	
}

\زیرمسئله[4]
یک الگوریتم  مناسب برای مرتب‌سازی $n$ عدد صحیح که بین $1$ تا $n^{100}$ هستند پیشنهاد کنید.

\پاسخ{
مرتب‌سازی مبنایی. همان‌طور که در کلاس گفته شد این مرتب‌سازی این لیست را در $O(n)$ مرتب می‌کند.
	
}
\نمره‌دهی{
	کلیت اثبات ۳ نمره و زمان اجرا یک نمره دارد. اگر فقط نام الگوریتم برده شده باشد ۱ نمره داده می‌شود. اگر الگوریتم غلط باشد نمره‌ای داده نمی‌شود.
	
}

\زیرمسئله[4]
فرض کنید می‌خواهیم کتاب‌های یک کتابخانه را بر حسب ارتفاعشان مرتب کنیم. ضخامت کتاب‌ها با هم مساوی است و برای هر مقایسه باید دو کتاب را از قفسه‌ها در بیاوریم و با هم مقایسه کنیم. الگوریتم پیشنهادی شما برای مرتب‌کردن کتاب‌های این کتابخانه چیست؟

\پاسخ{
مرتب‌سازی هرمی. زیرا هم زمانش در بدترین حالت $O(\log n)$ است و هم به حافظۀ اضافه (قفسه‌های خالی) نیاز ندارد. (مرتب‌سازی سریع نیز نمره می‌گیرد)
	
}
%heapsort - since it is in-place
\نمره‌دهی{
	کلیت الگوریتم ۲ نمره، زمان اجرا ۱ نمره و دلیل انتخاب ۱ نمره دارد.
	
}

\زیرمسئله[5]
الگوریتم پیشنهادی شما برای به دست آوردن مسیر بین دو رأس داده شده در یک درخت AVL  چیست؟

\پاسخ{
	مسیر هر رأس به ریشه را به دست آورده و در دو لیست می‌ریزیم که به خاطر متوازن بودن درخت AVL طول هر کدام $O(\log n)$ است. سپس از انتهای مسیر (ٰریشه) روی دو لیست حرکت می‌کنیم و قسمت مشترک را حذف می‌کنیم تا به پایین‌ترین جد مشترک برسیم. در انتها با برعکس کردن یک لیست و چسباندن این دو لیست به هم مسیر موردنظر پیدا می‌شود. زمان اجرای این الگوریتم $O(\log n)$ است.
	
}

\نمره‌دهی{
	مشابه ب.
	
}

\زیرمسئله[5]
برای بررسی این‌که یک درخت داده شده (که رئوس آن نیز رنگ شده‌اند) یک درخت قرمز-سیاه هست یا خیر یک الگوریتم پیشنهاد کنید.

\پاسخ{
ابتدا بررسی می‌کنیم که آیا درخت داده شده درخت دودویی جست‌وجو هست یا خیر. این کار را می‌توان با بررسی میان‌ترتیب رئوس انجام داد. سپس خاصیت‌های درخت قرمز سیاه را بررسی می‌کنیم. به غیر از خاصیت سیاه ارتفاع بقیه خاصیت‌ها را در هر رأس می‌توان بررسی کرد و سیاه ارتفاع را نیز با یک بار پیمایش درخت. در مجموع زمان اجرای الگوریتم $O(n)$ می‌باشد.
	
}

\زیرمسئله[5]
یک لیست به طول $n$ از اعداد متمایز داریم که از ابتدا تا یک خانه از آن (که نمی‌دانیم چه خانه‌ای است) صعودی و پس از آن نزولی است. یک الگوریتم کارا برای پیدا کردن مکان عدد داده شدۀ $x$ در این آرایه (و یا گزارش این‌که $x$ در این لیست وجود ندارد) پیشنهاد کنید.

\پاسخ{
	اگر بین هر دو خانه یک فلش فرضی درنظر بگیریم که به سمت خانۀ بزرگ‌تر اشاره می‌کند، می‌توان با جست‌وجوی دودویی روی این فلش‌ها خانۀ بیشینه را پیدا کرد. سپس لیست به دو قسمت صعودی و نزولی افراز می‌شود که $x$ را با جست‌وجوی دودویی در هر کدام جست‌وجو می‌کنیم.
	
}

\زیرمسئله[5]
یک واسط صف اولویت دو بعدی به این صورت تعریف می‌شود که کلید هر عنصر یک زوج مرتب $(x,y)$ است. علاوه بر درج 
(\lr{\texttt{insert(x, y, item)}})
 و حذف 
(\lr{\texttt{remove(x, y)}})
 یک عنصر با کلید مشخص، صف اولویت دو بعدی باید بتواند عنصر با کمترین $x$ (\متن‌تایپ{find-min-x()}) و عنصر با کمترین $y$ (\متن‌تایپ{find-min-y()}) را نیز برگرداند. داده‌ساختار پیشنهادی شما برای پیاده‌سازی صف اولویت دوبعدی چیست؟

\پاسخ{
راه‌حل این سؤال یکتا نیست. به طور مثال از دو درخت دودویی جست‌وجوی متوازن (مثل \متن‌لاتین{AVL}) استفاده می‌کنیم. در درخت اول کلیدها را ابتدا برحسب $x$ و اگر مساوی بودند برحسب $y$ مقایسه می‌کنیم. در درخت دوم کلیدها را ابتدا برحسب $y$ و اگر مساوی بودند برحسب $x$ مقایسه می‌کنیم. درج و حذف در دو درخت به طور عادی قابل انجام است. همچنین برای پیدا کردن عنصر با کمترین $x$ (\متن‌تایپ{find-min-x()}) از درخت اول و برای پیدا کردن عنصر با کمترین $y$ (\متن‌تایپ{find-min-y()}) از درخت دوم استفاده می‌کنیم. تمام عملیات در زمان $O(\log n)$ انجام می‌شود (راه‌حلی بهتری نیز وجود دارد که دو عملیات در $O(1)$ و دو تا در $O(\log n)$ انجام می‌شوند).
	
}

\نمره‌دهی{
اگر هر کدام از عملیات $\Theta(n)$ زمان بگیرد، ۱ نمره داده می‌شود. به طور مثال اگر به جای درخت AVL از هرم کمینه یا درخت دودویی جست‌وجوی عادی استفاده شود، \متن‌تایپ{remove} زمان خطی خواهد داشت.
	
}

\زیرمسئله[6]
یک مجموعه $n$تایی از نام‌های کاربری حداکثر ۲۰ حرفی داریم و می‌خواهیم ببینیم چند جفت نام کاربری وجود دارد که دقیقاً در یک حرف با هم تفاوت دارند. الگوریتمی از $O(n)$ طراحی کنید که این کار را انجام دهد. (توجه کنید در این مسئله از درهم‌سازی نمی‌توانید استفاده کنید)

\پاسخ{
	به ازای هر کدام از ۲۰ حرف مثل حرف $i$اُم، نام‌های کاربری را به ازای همۀ حروف به جز حرف $i$اُم با مرتب‌سازی مبنایی در $O(n)$ مرتب می‌کنیم (یعنی این الگوریتم ۲۰ بار مرتب‌سازی دارد). در این صورت رشته‌هایی که فقط در حرف $i$اُم تفاوت دارند پشت سر هم می‌آیند. با پیمایش لیست مرتب‌شده می‌توان این گروه‌های مشابه هم را پیدا کرد. اگر یک گروه با $k$ نام‌کاربری وجود داشته باشد که فقط در حرف $i$ با هم تفاوت دارند، می‌توان $\binom k 2$ را به جواب اضافه کرد. مجموع زمان اجرا $O(n)$ خواهد شد.
	
راه دیگر (که توسط دانشجویان پیشنهاد شده است): به جای مرتب‌سازی می‌توان نام‌کاربری‌ها را ۲۰ بار دوران داد به صورتی که هر بار حرف $i$اُم حرف آخر قرار بگیرد و آن‌ها را در ترای ریخت و با پیمایش ترای رشته‌های مشابه را پیدا کرد.
	
}

\نمره‌دهی{
	اگر (چه با استفاده از ترای و چه با استفاده از مرتب‌سازی) فقط اختلاف در حرف آخر را پیدا کنید، ۲ نمره داده می‌شود.
	
}

\زیرمسئله[6]
واسط دنبالۀ پویای محلی به این صورت تعریف می‌شود که علاوه بر لیست، یک مکان‌نما داریم که به یک عضو لیست اشاره می‌کند. عملیات دنبالۀ پویای محلی به صورت زیر است:
\شروع{فقرات}
\فقره
\متن‌تایپ{at(i)}: برگرداندن عنصر $i$اُم (برای $i$ دلخواه)
\فقره
\متن‌تایپ{cur-index()}: برگرداندن اندیس مکان‌نما.
\فقره
\متن‌تایپ{move-left()} و \متن‌تایپ{move-right()}: حرکت مکان‌نما یک واحد به چپ یا یک واحد به راست.
\فقره
\متن‌تایپ{delete-here()}: حذف عنصری که مکان‌نما روی آن است (و جابه‌جایی عناصر بعدی به چپ تا جای خالی پر شود).
\فقره
\متن‌تایپ{insert-here(x)} درج یک عنصر جدید در سمت چپ عنصری که مکان‌نما روی آن است (و جابه‌جا کردن عناصر بعدی) و حرکت مکان‌نما به روی عنصر جدید.
\پایان{فقرات}
چه داده‌ساختاری برای این واسط پیشنهاد می‌کنید؟

\پاسخ{
دو پشته یکی برای قسمتی از لیست که سمت راست مکان‌نما است و یکی برای قسمتی از لیست که سمت چپ مکان‌نما است می‌گیریم (به صورتی که عنصر مجاور مکان‌نما در بالای پشته قرار گیرد). همچنین عنصر جاری را در یک متغیر جداگانه ذخیره می‌کنیم. پشته‌ها را نیز با آرایۀ پویا پیاده‌سازی می‌کنیم 

برگرداندن عنصر جاری بدیهی است. اندیس مکان‌نما برابر با تعداد عناصر پشته چپ به اضافۀ یک است. حرکت به چپ و راست با یک pop از یک پشته و یک push به پشته دیگر قابل انجام است. همچنین حذف با یک pop و درج با یک push قابل انجام است. بنابراین تمام عملیات خواسته شده در $O(1)$ قابل انجام است.
	
}


\زیرمسئله[6]
در یک مهمانی $n$ مرد با بارانی‌های مشابه که فقط اندازه‌های متفاوت دارند وارد می‌شوند و آن‌ها را در یک اتاق قرار می‌دهند. در پایان مهمانی مشخص نیست کدام بارانی برای چه کسی هست. می‌خواهیم بارانی هر فرد را مشخص کنیم. تنها کاری که مجاز به انجام آن هستیم این است که یک بارانی را بر تن یک نفر امتحان کنیم که یا اندازه است، یا کوچک است و یا بزرگ. فرض کنید اندازۀ بارانی‌ها و افراد با هم متفاوت است و نمی‌توانیم دو بارانی و یا دو نفر را با هم مقایسه کنیم. چه الگوریتمی برای حل این مسئله پیشنهاد می‌کنید؟

\پاسخ{
	از الگوریتمی مشابه مرتب‌سازی سریع تصادفی استفاده می‌کنیم. یکی از افراد را به تصادف به عنوان لولا (محور) درنظر می‌گیریم و تمام بارانی‌ها را با او امتحان می‌کنیم. بنابراین یک بارانی اندازه، تعدادی کوچکتر و تعدادی بزرگتر می‌شود. سپس با استفاده از بارانی اندازه و امتحان کردن آن با همۀ افراد، افراد نیز به همین سه دسته افراز می‌شوند. فرد انتخابی و بارانی آن کنار گذاشته شده و دو دستۀ بارانی‌ها و افراد کوچکتر و بارانی‌ها و افراد بزرگتر را به صورت بازگشتی می‌توانیم حل کنیم. زمان اجرا به صورت تصادفی (میانگین) $O(n\log n)$ خواهد شد.
	
}

%\زیرمسئله[6]
%یک مجموعه از اعداد متمایز داریم و به دنبال سه تا از این اعداد هستیم که جمع آن‌ها برابر با صفر شود. الگوریتمی با زمان اجرای $O(n^2)$ طراحی کنید که یک سه‌تایی با این خاصیت (در صورت وجود) پیدا کند. 
%به طراحی الگوریتمی با زمان اجرای $O(n^2\log n)$ نیمی از نمره داده خواهد شد. (توجه کنید شما مجاز به استفاده از درهم‌سازی نیستید)
%
%\پاسخ{
%عناصر را در $O(n\log n)$ مرتب می‌کنیم. همۀ عناصر را به عنوان عنصر اول سه‌تایی امتحان می‌کنیم. سپس یک اشاره‌گر مربوط به عنصر دوم به کوچکترین عنصر و یک اشاره‌گر مربوط به عنصر سوم به بزرگترین عنصر قرار می‌دهیم. در هر مرحله اگر جمع این سه عنصر بزرگ‌تر از صفر بود، اشاره‌گر عنصر سوم را حرکت می‌دهیم و اگر بزرگ‌تر از صفر بود، اشاره‌گر عنصر دوم را حرکت می‌دهیم. همچنین به تکراری نبودن عناصر سه‌تایی توجه می‌کنیم. زمان اجرای این الگوریتم برای یک عنصر اول $O(n)$ و در مجموع $O(n^2)$ می‌شود.
%	
%}

\پایان{زیرمسائل}